% ============================================================
%  Experiment 7 --- Computer-Based Motor Speed Measurement
%  (IPC/desktop version, using the Advantech PCI-1750 DAQ card)
% ============================================================
\experimentheader{7}{Experiment 7}{Computer-Based Motor Speed Measurement}{Instrumentation Lab}

\begin{objectives}
  \begin{enumerate}
    \item To understand the principle and operation of computer-based measurement systems.
    \item To learn to acquire and process signals automatically using a data-acquisition (DAQ) system.
    \item To become familiar with signal conditioning for interfacing sensors with digital acquisition hardware.
    \item To measure and analyse the rotational speed of a DC motor using pulse signals.
    \item To develop skills in integrating hardware and software for real-time measurement and analysis.
  \end{enumerate}
\end{objectives}

\begin{equipment}
  \begin{itemize}
    \item Personal computer (PC) running Windows, with a C/C++ development environment
    \item Data Acquisition (DAQ) card: Advantech PCI-1750 (16 isolated digital inputs, 16 isolated digital outputs, 1 isolated counter, 1 timer, 2500\,VDC optical isolation), installed in the host PC, with driver and \textbf{Advantech Navigator.exe} installed
    \item Sensor trainer with DC motor and a 12-slot wheel-encoder pulse output
    \item Components for the student-designed filter and signal-conditioning circuit (op-amps, resistors, capacitors, Schmitt-trigger ICs, optocouplers)
    \item Digital storage oscilloscope (DSO), dual-channel, $\geq$50\,MHz, with passive probes
    \item Connecting wires and breadboard
  \end{itemize}
\end{equipment}

\begin{introduction}
  In most measurement processes it is often necessary to acquire data automatically, without a human operator. A computer-based instrument has several advantages over conventional measuring instruments: it can acquire measurements quickly and accurately, sustain real-time control, and operate for extended periods with consistent precision.\vspace*{10pt}

  A computer-based measurement system typically contains the following elements (Figure~\ref{fig:exp7-daq-overview}): a transducer, a signal-conditioning stage, data-acquisition hardware, and software running on the host PC.\newpage
\end{introduction}
\vspace*{1.5em}

\begin{boxedfigure}
  \centering
  \begin{tikzpicture}[node distance=8mm and 10mm]
    \node[block] (trans) {DC Motor\\Sensor Trainer};
    \node[block, right=of trans] (fil) {Filter\\Circuit};
    \node[block, right=of fil] (cond) {Signal\\Conditioning Circuit};
    \node[block, right=of cond] (daq) {PCI-1750\\DAQ Card};
    \node[block, below=of daq] (pc) {PC / C\\Software};

    \draw[arr] (trans) -- (fil);
    \draw[arr] (fil) -- (cond);
    \draw[arr] (cond) -- (daq);
    \draw[arr] (daq.south) -- (pc.north);
  \end{tikzpicture}
  \captionof{figure}{Computer-based measurement system: transducer, signal conditioning, DAQ hardware, and analysis software.}
  \label{fig:exp7-daq-overview}
\end{boxedfigure}
\vspace*{1.5em}

\begin{safetybox}
  \begin{enumerate}
    \item \textbf{Maximum voltage tolerance (0--5\,V TTL limit).} The PCI-1750's digital input channels require standard 0--5\,V TTL logic levels. Never connect a raw sensor signal exceeding 5\,V directly to the DAQ card input. Measure the raw pulse voltage amplitude with the DSO \emph{before} making any DAQ card connection.
  \end{enumerate}
\end{safetybox}

\begin{procedure}

  \textbf{Activity 1 --- Characterising the encoder signal and designing the filter}\vspace{2pt}

  \textit{Objective: measure the encoder's pulse characteristics and design a signal-conditioning filter.}

  \begin{enumerate}
    \item Power on the sensor trainer and verify motor rotation. Wire the encoder's raw output test point to Channel 1 of the oscilloscope.
    \item Run the motor across its usable speed range and observe the encoder pulse train on the oscilloscope. Record the peak-to-peak voltage, rise time, overshoot, and any noise ripple on the raw signal.
    \item At a representative operating speed, measure the pulse period and estimate the corresponding fundamental frequency from the oscilloscope trace.
    \item From this measurement, determine a suitable cutoff frequency for the noise-rejection filter.
    \item Design the necessary filter based on the cutoff frequency determined above. Justify your design choices.
    \item Design the signal-conditioning circuit needed to present the filtered encoder pulses to the PCI-1750's digital input with valid 0--5\,V TTL logic levels. Justify your design choices.
  \end{enumerate}
  \newpage

  \textbf{Activity 2 --- Interfacing to the DAQ hardware}\vspace{2pt}

  \textit{Objective: connect the conditioned encoder signal to the PC via the PCI-1750, and verify the connection before writing any software.}

  \begin{enumerate}
    \item With the motor trainer powered off, connect the output of your signal-conditioning circuit to the designated digital input channel (e.g.\ DI0 / Counter In) on the PCI-1750 terminal board. Connect a common ground between your breadboard and the PCI-1750 terminal board ground.
    \item On Channel 2 of the oscilloscope, verify that the conditioned signal is a clean 0--5\,V square wave with sharp edges and no false-trigger spikes, while keeping Channel 1 on the raw input for comparison. Adjust your filter's RC components if ringing occurs.
    \item Launch \textbf{Advantech Navigator.exe} on the PC to verify hardware channel detection and pulse-counting activity --- use this strictly as a hardware diagnostic to confirm the expected pulses are being registered, before writing any code.
    \item Close Navigator.exe and launch your C development environment.
  \end{enumerate}

  \begin{boxedfigure}
    \centering
    \resizebox{0.98\linewidth}{!}{%
      \begin{tikzpicture}[node distance=12mm and 16mm]
        \node[block, minimum width=2.6cm] (trainer) {DC Motor\\Trainer};
        \node[block, right=of trainer, minimum width=2.8cm] (filt) {Filter \&\\Signal Cond.\\ Circuit};
        \node[block, right=of filt, minimum width=2.6cm] (daq) {PCI-1750\\DAQ Card};
        \node[block, right=of daq, minimum width=2.4cm] (pc) {PC\\\scriptsize C Software};

        \draw[arr] (trainer) -- (filt);
        \draw[arr] (filt) -- (daq);
        \draw[arr] (daq) -- (pc);
      \end{tikzpicture}
    }
    \captionof{figure}{Signal path from the DC motor trainer to the host PC.}
    \label{fig:exp7-signalpath}
  \end{boxedfigure}

  \textbf{Activity 3 --- Speed measurement software}\vspace{2pt}

  \textit{Objective: implement software to compute and display motor speed (frequency and RPM) in real time, using two selectable measurement methods.}

  Two complementary methods can be used to convert the pulse train into a frequency and rotational speed, and your software should let the user select between them:
  \vspace*{1.5em}

  \textbf{Method 1 --- Fixed-Time (by counts, $1\leq N \leq 99$):} the user specifies a target pulse count $N$; the software measures the elapsed time $T_N$ for $N$ pulses to occur, then computes
  \[
    f = \frac{N}{T_N \times P} \ \left[\text{rev/s}\right], \qquad
    \text{RPM} = 60f, \qquad
    \omega = 2\pi f \ \left[\text{rad/s}\right]
  \]
  where $P$ is the encoder's pulse resolution (pulses per revolution).
  \vspace*{1.5em}

  \textbf{Method 2 --- Fixed-Count (by period, $1\leq T \leq 9$\,s):} the user specifies a measurement window $T$ seconds; the software counts the number of pulses $N_T$ that occur within that window, then computes
  \[
    f = \frac{N_T}{T \times P} \ \left[\text{rev/s}\right], \qquad
    \text{RPM} = 60f, \qquad
    \omega = 2\pi f \ \left[\text{rad/s}\right]
  \]
  \newpage

  In both methods, $t_{\text{start}}$ and $t_{\text{stop}}$ (marking the beginning and end of the measurement) are obtained from the PCI-1750's counter/timer using
  \[
    \text{Elapsed time} = \frac{t_{\text{stop}} - t_{\text{start}}}{\text{CLK\_TCK}}
  \]
  where $\text{CLK\_TCK}$ is the counter's clock tick rate.

  \begin{enumerate}
    \item In C, write a program that:
          \begin{itemize}
            \item initialises the PCI-1750 via the vendor driver/SDK (counter and digital-input registers).
            \item lets the user select the measurement method (Fixed-Time, $1 \leq N \leq 99$; or Fixed-Count, $1 \leq T \leq 9$\,s) and enter the corresponding parameter.
            \item for the \emph{Fixed-Time} method: times how long it takes for $N$ encoder pulses to occur.
            \item for the \emph{Fixed-Count} method: counts how many encoder pulses occur within the window $T$.
            \item computes the frequency, RPM, and angular speed $\omega$ using the relations above, updating at least every 500\,ms.
            \item displays the running value (frequency and RPM, with the active method indicated) in a simple HMI, updating in real time, with out-of-range error detection.
          \end{itemize}
    \item Interface your filter/conditioning circuit and the PCI-1750 together and perform measurements across the motor's input voltage range (e.g.\ 3.0\,V to 12.0\,V in 1.5\,V increments), using both methods at each operating point. Compare your computed speed against a direct oscilloscope measurement of the encoder period at the same operating point.
    \item Tabulate your results (Tables~\ref{tab:exp7-bycounts-rpm}--\ref{tab:exp7-bytime-freq}).
    \item Plot computed speed (from both software methods) against the oscilloscope-measured speed, and plot relative error (\%) against motor speed for both methods.
    \item Discuss and comment on your results, including any discrepancy between the two software methods and the oscilloscope reference, and how each method's resolution changes with motor speed (see the trade-off noted above).
  \end{enumerate}

  \needspace{6cm}
  \begin{boxedtable}
    \captionof{table}{By-counts method: frequency comparison table.}
    \label{tab:exp7-bycounts-freq}
    \centering
    {\renewcommand{\arraystretch}{1.4}
      \begin{tabular}{|c|c|c|c|c|c|c|}
        \hline
        \multirow{2}{*}{$N$} & \multicolumn{5}{c|}{Measured Frequency (Hz)} & \multirow{2}{*}{\shortstack{Oscilloscope\\(Hz)}}                     \\
        \cline{2-6}
                             & \#1                                          & \#2                                              & \#3 & \#4 & \#5 & \\
        \hline
                             &                                              &                                                  &     &     &     & \\
        \hline
                             &                                              &                                                  &     &     &     & \\
        \hline
                             &                                              &                                                  &     &     &     & \\
        \hline
      \end{tabular}
    }
  \end{boxedtable}
  \newpage

  \begin{boxedtable}
    \captionof{table}{By-counts method --- RPM comparison table.}
    \label{tab:exp7-bycounts-rpm}
    \centering
    {\renewcommand{\arraystretch}{1.4}
      \begin{tabular}{|c|c|c|c|c|c|c|}
        \hline
        \multirow{2}{*}{$N$} & \multicolumn{5}{c|}{$\omega$ (Measured RPM)} & \multirow{2}{*}{\shortstack{Oscilloscope\\(calculated RPM)}}                     \\
        \cline{2-6}
                             & \#1                                          & \#2                                                          & \#3 & \#4 & \#5 & \\
        \hline
                             &                                              &                                                              &     &     &     & \\
        \hline
                             &                                              &                                                              &     &     &     & \\
        \hline
                             &                                              &                                                              &     &     &     & \\
        \hline
      \end{tabular}
    }
  \end{boxedtable}
  \vspace*{10mm}

  \begin{boxedtable}
    \captionof{table}{By-period method: frequency comparison table.}
    \label{tab:exp7-bytime-freq}
    \centering
    {\renewcommand{\arraystretch}{1.4}
      \begin{tabular}{|c|c|c|c|c|c|c|}
        \hline
        \multirow{2}{*}{$T$ (s)} & \multicolumn{5}{c|}{Measured Frequency (Hz)} & \multirow{2}{*}{\shortstack{Oscilloscope\\(Hz)}}                     \\
        \cline{2-6}
                                 & \#1                                          & \#2                                              & \#3 & \#4 & \#5 & \\
        \hline
                                 &                                              &                                                  &     &     &     & \\
        \hline
                                 &                                              &                                                  &     &     &     & \\
        \hline
                                 &                                              &                                                  &     &     &     & \\
        \hline
      \end{tabular}
    }
  \end{boxedtable}
  \vspace*{10mm}

  \begin{boxedtable}
    \captionof{table}{By-period method: RPM comparison table.}
    \label{tab:exp7-bytime-rpm}
    \centering
    {\renewcommand{\arraystretch}{1.4}
      \begin{tabular}{|c|c|c|c|c|c|c|}
        \hline
        \multirow{2}{*}{$T$ (s)} & \multicolumn{5}{c|}{$\omega$ (Measured RPM)} & \multirow{2}{*}{\shortstack{Oscilloscope\\(calculated RPM)}}                     \\
        \cline{2-6}
                                 & \#1                                          & \#2                                                          & \#3 & \#4 & \#5 & \\
        \hline
                                 &                                              &                                                              &     &     &     & \\
        \hline
                                 &                                              &                                                              &     &     &     & \\
        \hline
                                 &                                              &                                                              &     &     &     & \\
        \hline
      \end{tabular}
    }
  \end{boxedtable}

\end{procedure}

\begin{reportq}
  \begin{enumerate}
    \item Present your oscilloscope measurements of the raw and conditioned encoder pulse, and your resulting filter design (topology, calculations, and component values), together with the measured response of the built filter.
    \item Present your design of the signal-conditioning stage between the filter and the PCI-1750 digital input, with justification, including the isolation and grounding measures you applied (see Safety \& Hardware Protection above).
    \item Compare your software-computed speed (both Fixed-Time and Fixed-Count methods) with the oscilloscope-measured speed across your tested operating points, and discuss the sources of any discrepancy.
    \item Compare the Fixed-Time and Fixed-Count methods directly: at what motor speeds does each method give better resolution or a more stable reading, and why? Which would you recommend for a low-speed application, and which for a high-speed one?
    \item What is the theoretical maximum measurable speed of this system, given the encoder resolution and your software's timing resolution?
    \item Identify the primary sources of measurement error in your system (e.g.\ pulse jitter, counter quantisation error, sampling clock drift, motor speed ripple), and discuss their impact on overall RPM accuracy.
  \end{enumerate}
\end{reportq}

\begin{labsection}{Deliverables}{}
  \begin{enumerate}
    \item \textbf{C source code file (.c):} fully commented C software implementing pulse counting, timing calculations, RPM conversion, and HMI display.
    \item \textbf{Technical design report (.pdf):} formal lab report detailing circuit schematics, component calculations, signal waveform captures, data tables, and error analysis.
    \item \textbf{System demonstration:} a live bench demonstration showing real-time motor speed tracking on your C software HMI, alongside DSO verification.
  \end{enumerate}
\end{labsection}

\begin{labsection}{References}{}
  \begin{itemize}
    \item Advantech PCI-1750/PCI-1750SO user manual.
  \end{itemize}
\end{labsection}
